% ============================================================
%  05_introduction_generale.tex  –  Introduction Générale
% ============================================================
\chapter*{Introduction Générale}
\addcontentsline{toc}{chapter}{Introduction Générale}
\markboth{Introduction Générale}{Introduction Générale}

La transformation numérique a profondément modifié les pratiques d'apprentissage,
aussi bien dans les établissements d'enseignement que dans les entreprises. Dans un
contexte marqué par la formation à distance, la montée en puissance de la formation
continue et l'évolution rapide des compétences numériques, les plateformes e-learning
sont devenues des outils essentiels pour diffuser, suivre et évaluer les apprentissages.
\vspace{1cm}
\medskip


Toutefois, de nombreuses solutions existantes restent limitées par une faible personnalisation
des parcours, des évaluations essentiellement quantitatives et une assistance pédagogique
insuffisante. L'émergence de l'intelligence artificielle générative, notamment des grands
modèles de langage, offre de nouvelles perspectives pour accompagner les apprenants,
automatiser certaines tâches pédagogiques et proposer des recommandations adaptées
aux besoins réels des étudiants.
\vspace{1cm}

\medskip

C'est dans ce cadre que s'inscrit le présent projet de fin d'études, réalisé au sein de
Whitecape Technologies. Il porte sur la conception et le développement d'\textbf{EduBoard},
une plateforme e-learning intelligente visant à combiner gestion pédagogique, suivi
de progression, évaluation assistée par \ac{IA}, recommandation personnalisée, chatbot
pédagogique et génération automatique de certificats.
\vspace{1cm}

\medskip

Ce rapport est organisé en quatre chapitres principaux, encadrés par une introduction
et une conclusion générales~:
\vspace{0.5cm}

\begin{itemize}
  \item \textbf{Chapitre~1} présente le contexte général du projet~: l'organisme
        d'accueil, la problématique, l'étude de l'existant, la solution proposée et la
        méthodologie de travail.
  \item \textbf{Chapitre~2} expose la spécification détaillée des besoins fonctionnels
        et non fonctionnels, avec les diagrammes \ac{UML} de cas d'utilisation et les
        scénarios principaux.
  \item \textbf{Chapitre~3} détaille la conception globale de la solution~: architecture,
        modèle de données, diagrammes d'activité, de séquence, de classes et de
        déploiement.
  \item \textbf{Chapitre~4} présente la réalisation concrète, les interfaces développées
        et les tests de validation effectués.
\end{itemize}
